\begin{center}
{\Large \textbf{Data Management Plan}}
\end{center}
\vspace*{-8pt}

\paragraph{Types of Data}

This project will generate and use five categories of data. \textbf{Policy rule corpora}: structured constraint schemas extracted from real and synthetic organizational policy documents (scholarship handbooks, healthcare compliance policies, government guidelines), plus expert-curated gold-standard rule sets for evaluation. Estimated volume: $<$5\,MB of structured JSON/YAML. \textbf{SimBench v2 benchmark}: approximately 130 human-validated evaluation tasks across three policy-constrained domains (scholarship and financial aid, regulated customer service, enterprise policy compliance), each with ground-truth governance outcomes, domain-expert validations, and task metadata. Estimated volume: $<$50\,MB. \textbf{Adversarial test suite}: 600 labeled test cases (200 per attack vector: preference injection, authority confusion, multi-step accumulation), each with expected governance outcome, attack label, and model family. Estimated volume: $<$20\,MB. \textbf{Resolution trace logs}: structured JSON decision-governance records from evaluation runs on SimBench v2, the adversarial suite, and $\tau$-bench. These are synthetic or de-identified; no live organizational data is retained. Estimated volume: $<$1\,GB. \textbf{Compliance reviewer study data} (Aim~3): within-subjects human-subjects data including task accuracy judgments, time-to-determination, reviewer confidence ratings, and de-identified demographic information. All interactions use synthetic or de-identified scenarios from WSOSF; no real student records are used. Estimated volume: $<$1\,MB.

\paragraph{Metadata Standards}

All datasets are documented with standardized metadata following FAIR principles~\cite{wilkinson2016fair}: Title, Creator, Date, Format, Subject, Persistent Identifier, Privacy Classification, and Description. Dataset-specific fields include Collection Phase, IRB Protocol Reference (Aim~3), Validation Method (expert or crowd), Policy Domain, Attack Vector (adversarial suite), and Retention Period. Policy corpora and benchmark tasks use JSON-Schema-documented formats. Resolution traces follow the project's open specification (released Y3~Q4). Human-subjects data include codebooks for all instruments and anonymization documentation.

\paragraph{Policies for Access, Sharing, and Privacy}

The project implements a tiered access model. \textbf{Public (Tier~1)}: SimBench v2 (CC-BY~4.0 with online leaderboard), the 600-case adversarial test suite, the resolution trace format specification, and the governance layer SDK (open-source). \textbf{Restricted (Tier~2)}: de-identified Aim~3 study data (reviewer accuracy, timing, confidence) available to approved researchers under a data use agreement. \textbf{Confidential (Tier~3)}: linking keys, raw WSOSF policy documents provided under confidentiality agreement. No personal interaction data from live users is collected; all Aim~3 scenarios are synthetic or de-identified before reaching the research team. WSOSF policy materials are governed by a data-sharing agreement specifying retention, destruction, and re-use limitations. Aim~3 human-subjects activities will receive IRB approval (anticipated Exempt or Expedited review) prior to recruitment; the IRB protocol is described in~\S2.3.

\paragraph{Policies for Re-use, Re-distribution, and Derivatives}

SimBench v2 and the adversarial test suite are released under CC-BY~4.0, permitting re-use, redistribution, and derivative works with attribution. The resolution trace format specification is released under CC-BY~4.0; contributions to standardization bodies (IEEE, ISO AI governance working groups) will be governed by those bodies' intellectual property policies. The governance layer SDK is released under the Apache 2.0 license. Aim~3 de-identified data may be redistributed to researchers under a data use agreement; no re-identification or linkage to WSOSF individuals is permitted by any downstream user. The PI will publish sufficient methodological detail to support independent replication of key findings.

\paragraph{Plans for Archiving and Preservation}

Datasets and code will be deposited in the UW ResearchWorks Archive (institutional DOI) and mirrored in a public GitHub organization. Deposits are made at each major milestone (see Table~1 in the Project Description). Data are retained for a minimum of three years following project completion, or longer per IRB protocol or journal policy. Aim~3 linking keys are destroyed upon completion of data collection and final dataset validation. Open-source repositories are maintained with tagged releases corresponding to published results. Version control (Git) is used throughout the project for all code and data generation pipelines.

\paragraph{Roles and Responsibilities}

The PI (Dr.\ Shah) maintains overall data management responsibility. The postdoctoral researcher (Years~2--3) oversees day-to-day Aim~3 data collection, IRB compliance, and dataset documentation. The GRA (Harshita Chopra) manages evaluation infrastructure, benchmark construction pipelines, and WSOSF scenario collection. All team members complete required CITI human-subjects research training prior to accessing Aim~3 data. UW Information Technology provides encrypted storage and backup infrastructure.

\vspace*{-8pt}
\begingroup
\renewcommand{\bibsection}{}
\bibliographystyle{plain}
\begin{thebibliography}{9}
\bibitem{wilkinson2016fair}
M.~D.~Wilkinson et al.
\newblock The {FAIR} guiding principles for scientific data management and stewardship.
\newblock \textit{Scientific Data}, 3:160018, 2016.
\end{thebibliography}
\endgroup
